%!TEX program = xelatex
\documentclass[12pt,a4paper]{exam-zh}

\usepackage{amsmath,array,graphicx,booktabs,multirow}
\usepackage{enumitem}
\usepackage{fancyhdr}
\usepackage{tikz}
\usepackage{tabularx}
\usepackage{setspace}
\setCJKmainfont{Noto Serif CJK SC}
\geometry{margin=2.2cm,top=2.0cm,bottom=2.2cm}
\pagestyle{fancy}
\fancyhf{}
\cfoot{-\thepage{}-}
\renewcommand{\headrulewidth}{0pt}
\renewcommand{\footrulewidth}{0pt}
\setlength{\parindent}{0pt}
\setlength{\parskip}{0.2em}
\setlength{\headheight}{14pt}
\linespread{1.12}
\usepackage{setspace}
\usepackage{eso-pic}
\usetikzlibrary{calc}

%region exam setup

\examsetup{
    page = {
    size = a4paper,           % 设定纸张大小 [4]
    foot-content = { - ; - }  % 页脚显示为 -1-, -2- 格式 [3]
  },
    choices = {
    label-pos = top-left,
    column-sep = 1em,
    label-pos = auto,
    label-sep = 0.5em,
    max-columns = 4
    },
  question = {
    top-sep    = 1em, % 调整题目上方的垂直距离（默认值为 .5em plus .5em minus .2em）
    bottom-sep = 1em  % 调整题目下方的垂直距离（默认值为 .5em plus .5em minus .2em）
  },
    paren = {
    show-paren = false  % 设置为 false 则不显示括号 [3, 4]
  },
    choices = {
    label = \textbf{\Alph*}
  }, 
    solution = {
    show-solution = show-move, % 答案自动移动到卷末
    label-content = {参考答案},  % 答案区的标题文字
    show-qed = true            % 在解答末尾显示结束符号 ◻ 
  }
}

%endregion


%region 卷面变量设置
    % ==========================
    % ==========================
% 卷面变量设置（只需改这里）
% ==========================
\newcommand{\papercode}{SC 07}
\newcommand{\examtitle}{2026年金宝培元独立中学半年考}
\newcommand{\examlevel}{高中组}
\newcommand{\examsubject}{高级数学(II)}
\newcommand{\sidebartext}{高级数学(II)/二}
\newcommand{\papername}{试卷一}
\newcommand{\papertype}{选择题}
\newcommand{\examdate}{2026年5月20日}
\newcommand{\examtime}{8:00 a.m. 至9:00 a.m.}
\newcommand{\examduration}{1小时}
\newcommand{\questionpages}{5面}
\newcommand{\paperonepercentage}{40\%}
\newcommand{\papertwopercentage}{60\%}
\newcommand{\setterteacher}{蔡杰庆师}
\newcommand{\reviewteacher}{王萌师}
\newcommand{\auditteacher}{雷泳健师}
    \newcommand{\optrow}[4]{%
    \begin{tabular*}{\linewidth}{@{\extracolsep{\fill}}llll}
    A~#1 & B~#2 & C~#3 & D~#4
    \end{tabular*}%
    }
    \newcommand{\qitem}[5]{\item #1\par\optrow{#2}{#3}{#4}{#5}}
%endregion
    \begin{document}


%region 封面
    \AddToShipoutPictureFG*{%
    \put(31,104){\color{black}\rule{16pt}{66pt}}%
    \put(37,112){\makebox(0,0)[lb]{\rotatebox{90}{\color{white}\bfseries\footnotesize \sidebartext}}}%
    \put(24,18){\makebox(0,0)[lb]{\small 版权所有}}%
    \put(520,18){\makebox(0,0)[lb]{\bfseries\small 【机密】}}%
    }

    \thispagestyle{fancy}

    \begin{minipage}[t]{0.22\linewidth}
    \vspace{0pt}
    \end{minipage}
    \hfill
    \begin{minipage}[t]{0.56\linewidth}
    \centering
    {\zihao{3}\bfseries \examtitle\par}
    \vspace{2.3em}
    {\zihao{3}\bfseries \examlevel\par}
    \vspace{0.3em}
    {\zihao{2}\bfseries \examsubject\par}
    \vspace{0.2em}
    {\large (\papercode)\par}
    \end{minipage}
    \hfill
    \begin{minipage}[t]{0.18\linewidth}
    \vspace{0pt}
    \end{minipage}

    \vspace{3.0em}

    \begin{center}
    \begin{tabular}{c}
        \papername\hspace{2em}\papertype\\[-0.1em]
        \midrule
        日期:\examdate\\[-0.1em]
        时间:\examtime\\[-0.1em]
        (\examduration)\\
        \midrule
    \end{tabular}
    \end{center}

    \vspace{2.0em}

    {\zihao{4}\bfseries \underline{考生须知}\par}
    \vspace{0.8em}

    \begin{enumerate}[label=\arabic*、,leftmargin=2em,itemsep=0.65em]
    \item 本科试卷共分两份：\\
    试卷一：选择题（\paperonepercentage）；\\
    试卷二：作答题（\papertwopercentage）。

    \item 考生须于第一阶段规定的60分钟内完成试卷一。第二阶段规定的2小时内作答试卷二。

    \item 试卷一选择题20题全答。选出正确的答案，然后用\textbf{5B铅笔将“O”答案纸}（电脑卡）上相应的字母所在的小方格涂黑。

    \item 可使用电子计算机进行演算。

    \item 试卷一选择题共印\questionpages。
    \end{enumerate}

    \vspace{2.0em}

    \begin{center}
    {\setlength{\fboxrule}{1.2pt}\fbox{\rule{0pt}{1.1em}\hspace{1.2em}{\zihao{4}\bfseries 未经正式宣布\quad 不得翻看内页}\hspace{1.2em}}}
    \end{center}

    \vfill

    \begin{tabular*}{\linewidth}{@{\extracolsep{\fill}}ll}
    考生姓名：\underline{\hspace{4.2cm}} & 考生编号：\underline{\hspace{4.0cm}}
    \end{tabular*}

%endregion

\newpage





\vfill % 与上方内容保持距离

\noindent
\newlength{\sigwidth}
\settowidth{\sigwidth}{（\hspace{4em}师）}

\begin{tabular*}{\textwidth}{@{\extracolsep{\fill}} l l l}
    出题： & 审题： & 审核： \\[4em]
    \rule{\sigwidth}{0.5pt} &
    \rule{\sigwidth}{0.5pt} &
    \rule{\sigwidth}{0.5pt} \\
    （\hspace{0em} \setterteacher ） &
    （\hspace{0em} \reviewteacher ）  &
    （\hspace{0em} \auditteacher ） \\
\end{tabular*}

\clearpage
\thispagestyle{empty} % 取消这一页的页眉页脚（页码）
~
\clearpage
\end{document}
